\documentclass[11pt]{article}
\usepackage[final]{acl}
\usepackage{times}
\usepackage{latexsym}
\usepackage[T1]{fontenc}
\usepackage[utf8]{inputenc}
\usepackage{microtype}
\usepackage{inconsolata}
\usepackage{booktabs}
\usepackage{graphicx}

% Generated by scripts/make_tables.py; do not edit.
\newcommand{\relSources}{41}
\newcommand{\relNational}{20}
\newcommand{\relNoCorpus}{12}
\newcommand{\relLangs}{224}
\newcommand{\relNoncommercial}{39}
\newcommand{\relArchiveGB}{1.26}
\newcommand{\relArchiveMedianKB}{107}
\newcommand{\relArchiveMaxMB}{372}
\newcommand{\relScored}{212}
\newcommand{\accOffN}{103}
\newcommand{\accOffMacro}{89.85}
\newcommand{\accOffCopy}{58.41}
\newcommand{\accOffPooled}{94.71}
\newcommand{\accOffTokens}{2,176,260}
\newcommand{\accTwoN}{109}
\newcommand{\accTwoMacro}{73.87}
\newcommand{\accTwoCopy}{57.54}
\newcommand{\accScoredOfficial}{103}
\newcommand{\speedRsOneMedian}{54.4M}
\newcommand{\speedRsOneMin}{6.2M}
\newcommand{\speedRsOneMax}{858M}
\newcommand{\speedRsSixteenMedian}{759M}
\newcommand{\speedRsSixteenMin}{93.2M}
\newcommand{\speedRsSixteenMax}{12708M}
\newcommand{\speedPyOneMedian}{18.0M}
\newcommand{\speedPyOneMin}{4.8M}
\newcommand{\speedPyOneMax}{28.0M}
\newcommand{\speedPySixteenMedian}{142M}
\newcommand{\loadRsMedianMs}{0.8}
\newcommand{\loadRsMaxMs}{1183}
\newcommand{\speedLangs}{169}
\newcommand{\speedChecked}{1,011,049,401}
\newcommand{\speedMismatch}{0}
\newcommand{\speedRivalLangs}{91}
\newcommand{\speedRsFastestLangs}{90}
\newcommand{\speedRsSlowerLangs}{kat}
\newcommand{\cellsNew}{2,636}
\newcommand{\cellsReused}{1,675}
\newcommand{\foreignNew}{0.08}
\newcommand{\foreignReused}{1.07}
\newcommand{\accsimplemmaN}{33}
\newcommand{\accsimplemmaTheirs}{90.75}
\newcommand{\accsimplemmaOurs}{95.44}
\newcommand{\acclemmagenThreeN}{16}
\newcommand{\acclemmagenThreeTheirs}{82.27}
\newcommand{\acclemmagenThreeOurs}{96.33}
\newcommand{\accudpipeOneN}{44}
\newcommand{\accudpipeOneTheirs}{91.36}
\newcommand{\accudpipeOneOurs}{95.16}
\newcommand{\accspacysmN}{17}
\newcommand{\accspacysmTheirs}{89.55}
\newcommand{\accspacysmOurs}{95.77}
\newcommand{\accspacylgN}{17}
\newcommand{\accspacylgTheirs}{90.63}
\newcommand{\accspacylgOurs}{95.77}
\newcommand{\acccorenlpN}{1}
\newcommand{\acccorenlpTheirs}{96.37}
\newcommand{\acccorenlpOurs}{97.87}
\newcommand{\accstanzaN}{65}
\newcommand{\accstanzaTheirs}{95.32}
\newcommand{\accstanzaOurs}{93.81}
\newcommand{\acctrankitN}{46}
\newcommand{\acctrankitTheirs}{91.32}
\newcommand{\acctrankitOurs}{94.44}
\newcommand{\accspacytrfN}{7}
\newcommand{\accspacytrfTheirs}{93.98}
\newcommand{\accspacytrfOurs}{97.14}
\newcommand{\accudpipeTwoN}{77}
\newcommand{\accudpipeTwoTheirs}{94.48}
\newcommand{\accudpipeTwoOurs}{93.47}
\newcommand{\speedsimplemmaN}{37}
\newcommand{\speedsimplemmaUp}{5.4}
\newcommand{\speedlemmagenThreeN}{17}
\newcommand{\speedlemmagenThreeUp}{9.5}
\newcommand{\speedudpipeOneN}{50}
\newcommand{\speedudpipeOneUp}{2,177}
\newcommand{\speedspacysmN}{20}
\newcommand{\speedspacysmUp}{1,422}
\newcommand{\speedspacylgN}{20}
\newcommand{\speedspacylgUp}{1,568}
\newcommand{\speedcorenlpN}{1}
\newcommand{\speedcorenlpUp}{1,965}
\newcommand{\speedstanzaGpu}{16k}
\newcommand{\speedstanzaGpuOneUp}{1,518}
\newcommand{\speedstanzaN}{90}
\newcommand{\speedstanzaUp}{8,160}
\newcommand{\speedtrankitGpu}{3.9k}
\newcommand{\speedtrankitGpuOneUp}{8,160}
\newcommand{\speedtrankitN}{53}
\newcommand{\speedtrankitUp}{433,457}
\newcommand{\speedspacytrfGpu}{20k}
\newcommand{\speedspacytrfGpuOneUp}{1,859}
\newcommand{\speedspacytrfN}{8}
\newcommand{\speedspacytrfUp}{143,271}
\newcommand{\accuracyRows}{simplemma & 33 & 90.75 & 95.44 & +4.68 & 33 / 0 \\
LemmaGen3 & 16 & 82.27 & 96.33 & +14.06 & 16 / 0 \\
UDPipe 1 & 44 & 91.36 & 95.16 & +3.81 & 38 / 6 \\
spaCy sm & 17 & 89.55 & 95.77 & +6.22 & 15 / 2 \\
spaCy lg & 17 & 90.63 & 95.77 & +5.14 & 14 / 3 \\
CoreNLP & 1 & 96.37 & 97.87 & +1.50 & 1 / 0 \\
Stanza & 65 & 95.32 & 93.81 & -1.51 & 10 / 55 \\
Trankit & 46 & 91.32 & 94.44 & +3.12 & 29 / 17 \\
spaCy trf & 7 & 93.98 & 97.14 & +3.16 & 6 / 1 \\
UDPipe 2 & 77 & 94.48 & 93.47 & -1.01 & 12 / 63 \\
}
\newcommand{\speedRows}{simplemma & 37 & 6.6M & 5.4 & 2.2 & 57.8M & 9.2 & 324 / 110 \\
LemmaGen3 & 17 & 3.5M & 9.5 & 3.9 & 32.5M & 15.8 & 1,445 / 118 \\
UDPipe 1 & 50 & 15k & 2,177 & 910 & 151k & 3,481 & 153 / 99 \\
spaCy sm & 20 & 27k & 1,422 & 586 & 218k & 2,624 & 773 / 114 \\
spaCy lg & 20 & 25k & 1,568 & 646 & 165k & 3,304 & 1,546 / 114 \\
CoreNLP & 1 & 27k & 1,965 & 696 & 336k & 2,321 & 1,600 / 111 \\
Stanza & 90 & 4.5k & 8,160 & 3,447 & 16k (GPU) & 22,439 & 1,240 / 99 \\
Trankit & 53 & 74 & 433,457 & 186,480 & 3.9k (GPU) & 121,301 & 4,665 / 94 \\
spaCy trf & 8 & 278 & 143,271 & 55,319 & 20k (GPU) & 28,092 & 2,840 / 111 \\
}


\title{udlex: Dictionary Lemmatization for \relLangs{} Languages\\
at Tens of Millions of Words per Second on One Core}

\author{Maksim Kleshchenok \\
  Independent Researcher}

\begin{document}
\maketitle

\begin{abstract}
We present udlex, a lemmatizer that answers from compiled dictionaries and
falls back to learned suffix rules, with no neural model at inference. One
command rebuilds its dictionaries for \relLangs{} languages from Universal
Dependencies 2.19 and \relSources{} pinned lexical sources, including
Wiktionary, UniMorph and \relNational{} national resources. On the official
test splits of \accOffN{} languages it lemmatizes \accOffMacro\% of gold tokens
correctly (macro average), against \accOffCopy\% for copying the form. A
Rust core with Python, C, Go and OpenSearch interfaces processes a median of
\speedRsOneMedian{} words per second on one CPU core across \speedLangs{}
languages, and is the fastest of ten measured systems on \speedRsFastestLangs{}
of the \speedRivalLangs{} languages where another system was timed. Neural
toolkits remain more accurate where they have models: Stanza reaches
\accstanzaTheirs\% against \accstanzaOurs\% on \accstanzaN{} shared languages,
at a throughput thousands of times lower. All dictionaries, receipts and scripts
are public.
\end{abstract}

\section{Introduction}

Lemmatization is a preprocessing step in search engines, keyword extraction
and feature pipelines for named entity recognition, where a corpus may hold
billions of tokens and the step runs on every document and every query.
Neural toolkits such as Stanza \citep{qi-etal-2020-stanza}, Trankit
\citep{nguyen-etal-2021-trankit} and UDPipe~2 \citep{straka-2018-udpipe} are
the most accurate option for the languages they cover, but they tag
part of speech first, need a model per language and process at most tens
of thousands of words per second on one CPU core. Lookup lemmatizers such as simplemma
\citep{barbaresi-simplemma} and the lookup tables of spaCy
\citep{honnibal-2020-spacy} are fast, yet each covers a few dozen
languages with one resource per language.

udlex sits between the two. For each language it compiles a corpus table
learned from a Universal Dependencies treebank \citep{nivre-etal-2020-universal,ud219},
one or more external morphological dictionaries in a learned order of
priority, and a set of suffix rules for words that no table knows. Answers
are deterministic and every answer reports the stage that produced it. The
system makes three contributions:
\begin{itemize}
\item \textbf{Dictionaries for \relLangs{} languages}, rebuilt from pinned
inputs by one command and packaged with a license for each of their parts;
\relNoncommercial{} packages contain noncommercial data and are marked as
such (\S\ref{sec:dictionaries}).
\item \textbf{A compiled lookup engine} in Rust, memory-mapped, validated
when opened, and reachable from Python, C, Go and an OpenSearch analysis
plugin (\S\ref{sec:engine}).
\item \textbf{An evaluation on every language} with a test split, against
nine other lemmatizers on identical test words, with speed and memory
measured for all systems on one machine (\S\ref{sec:evaluation}).
\end{itemize}

\section{The lookup cascade}
\label{sec:engine}

udlex takes a sentence as a list of tokens and returns one lemma per token.
Each token goes through a fixed cascade of stages and the first stage that
knows the token answers:
\begin{enumerate}
\item \emph{entities}: curated proper names, first with the exact key, then
lowercased;
\item \emph{exact}: the corpus table, keyed by the surface form;
\item \emph{lower}: the lowercased corpus table, and the exact table retried
with a lowercased key;
\item \emph{external}: the external dictionaries, exact and then lowercased;
\item \emph{rule}: the longest matching learned suffix rewrite;
\item \emph{identity}: the form itself.
\end{enumerate}
The external table holds the dictionaries merged in a learned priority
order: when two sources disagree on a form, the earlier source owns the
entry. Per-language options, such as an adapter that maps dictionary lemmas to the
treebank's conventions or a compound-boundary marker, are decided on
development data when the dictionary is built and stored with it.

Tables are compiled into a bundle that the engine memory-maps. A table is
stored as a double-array trie, a finite-state transducer, or a hybrid that
keeps keys of up to eight bytes in a double array and longer keys in a hash
index with the lemma inline. All tables of the release that hold external
dictionaries use the hybrid, the fastest of the codecs on these lookups in
our measurements. A lemma that equals
its form is not copied: the engine returns a borrowed slice of the input, so
the common path allocates nothing. Bundles carry checksums and are validated
when opened, so a truncated or corrupted file fails to load instead of
returning wrong lemmas.

\paragraph{Interfaces.} The Python package wraps the Rust core through PyO3:
{\small
\begin{verbatim}
import udlex
lex = udlex.load("bundles", "fin")
lex.lemmatize(["Kirjakaupat", "olivat"])
lex.lemmatize_many(sentences)
lex.lemmatize_with_info(tokens)  # + stage
\end{verbatim}
}
A command-line tool lists, compiles and verifies bundles and maintains a
user dictionary, a small overlay of corrections with atomic updates and
rollback that takes precedence over the shipped tables. The same core is
exposed as a C library, through Go bindings, and as an OpenSearch 3.10
analysis plugin that ships with all \relLangs{} dictionaries.

\section{Building the dictionaries}
\label{sec:dictionaries}

A single script fetches every input by URL and SHA-256, splits the treebanks,
extracts form--lemma pairs, decides the source order, packages and evaluates.
Each package records the hash of every input it was built from.

\paragraph{Corpora.} Each language anchors on one UD~2.19 treebank. Of the
\relLangs{} languages, \accOffN{} have an official train and test split. For
\accTwoN{} test-only treebanks the evaluation uses two folds: a dictionary
learned on one half is scored on the other, and the shipped dictionary is
learned on all of it. \relNoCorpus{} languages have no treebank in UD~2.19 and
ship external dictionaries only. When a treebank lacks a development split, a
tenth of the training documents, chosen by hash, serves as one.

\paragraph{External dictionaries.} The \relSources{} sources are Wiktionary
editions extracted with Wiktextract \citep{ylonen-2022-wiktextract},
UniMorph \citep{kirov-etal-2018-unimorph}, MULTEXT-East
\citep{erjavec-2012-multext}, CC-CEDICT and \relNational{} national
resources, among them the Polish SGJP and PoliMorf, Slovenian Sloleks,
Icelandic B\'IN, Swedish SALDO, Latvian T\=ezaurs and Russian OpenCorpora.
A license policy excludes sources with copyleft or unclear terms; the
packages are under CC~BY-SA, CC~BY or CC~BY-NC-SA.

\paragraph{Source order.} The external dictionaries start in a declared order
by family: national resources and CC-CEDICT, then UniMorph, then Wiktionary. An order
learned on held-out training words replaces it only when it fixes
significantly more development words than it breaks (McNemar test,
$p \le 0.05$, with a minimum net gain); per-language options pass the same
gate. Without development data the declared order stays.

\paragraph{Licenses.} A package records the license of each table and the
attribution of each source, in the format of spaCy's package metadata. A
package containing any noncommercial part is named \texttt{-nc}. The
\relLangs{} packages take \relArchiveGB{}~GB compressed (median
\relArchiveMedianKB{}~kB, largest \relArchiveMaxMB{}~MB for Finnish) and are
distributed as release assets, because the largest exceed the size limit of
the Python package index.

\section{Evaluation}
\label{sec:evaluation}

\subsection{Accuracy}

We score exact lemma match on gold tokens of the UD~2.19 test split. On the
\accOffN{} languages with an official split, udlex reaches a macro average of
\accOffMacro\% (\accOffPooled\% of all \accOffTokens{} tokens), against
\accOffCopy\% for copying the form. On the \accTwoN{} two-fold languages it
reaches \accTwoMacro\%, against \accTwoCopy\%; these treebanks are small,
often a few hundred sentences.

Table~\ref{tab:accuracy} compares udlex with other lemmatizers on identical
test words. Each row uses only the languages that system covers, so rows are
not comparable with each other. We drop a language for a system when its
model was trained on data containing more than 1\% of our test sentences
(the UD release behind each model is known), when its Stanza package has no
lemmatizer, or when Stanza's lemma dictionary already knows at least half of
the test words that our training split lacks.

\begin{table}[t]
\centering\small
\setlength{\tabcolsep}{2.8pt}
\begin{tabular}{lrrrrr}
\toprule
System & Langs & Theirs & udlex & $\Delta$ & udlex $>$ / $<$ \\
\midrule
\accuracyRows
\bottomrule
\end{tabular}
\caption{Macro-averaged lemma accuracy (\%) on identical UD~2.19 test words.
Each row averages over the languages that system covers after removing test
leakage; the last column counts languages where udlex is more / less
accurate.}
\label{tab:accuracy}
\end{table}

udlex is more accurate than every dictionary or statistical lemmatizer on its
shared languages: \accudpipeOneOurs\% against \accudpipeOneTheirs\% for
UDPipe~1 on \accudpipeOneN{} languages and \accsimplemmaOurs\% against
\accsimplemmaTheirs\% for simplemma on \accsimplemmaN. Against Trankit
it is ahead on average but not everywhere. The two neural toolkits trained on
recent UD releases remain more accurate: Stanza reaches \accstanzaTheirs\%
against \accstanzaOurs\% for udlex on \accstanzaN{} languages, and UDPipe~2
\accudpipeTwoTheirs\% against \accudpipeTwoOurs\% on \accudpipeTwoN. They see the sentence and
the part of speech; udlex sees one token at a time.

\subsection{Speed and memory}

\begin{table*}[t]
\centering\small
\begin{tabular}{lrrrrrrr}
\toprule
& & \multicolumn{3}{c}{One core} & \multicolumn{2}{c}{Parallel} & Peak RSS, MB \\
\cmidrule(lr){3-5}\cmidrule(lr){6-7}
System & Langs & words/s & udlex-rs $\times$ & udlex-py $\times$ & words/s & udlex-rs $\times$ & theirs / udlex-rs \\
\midrule
\speedRows
\bottomrule
\end{tabular}
\caption{Throughput of other systems and udlex's speed-up over them, on the
languages each system was timed on. Rates are medians over languages;
speed-ups are geometric means of per-language ratios. ``Parallel'' is the
better of 16 processes and 16 threads on 16 physical cores, or one RTX~4090
for neural systems, against udlex's Rust core on 16 threads.}
\label{tab:speed}
\end{table*}

All systems were timed on one machine, an AMD Ryzen~9 7950X with 16 cores
and an RTX~4090, on the UD~2.19 test sentences of all \speedLangs{} languages
that have both a test text and a dictionary. A cell is one language, system
and mode in a fresh process pinned to its cores: one core, 16 processes or
16 threads on one hardware thread of each physical core, or the GPU. A cell
times three windows of at least a second, each on a new sample of sentences
drawn with replacement; we report the median over three rounds in shuffled
order. Input is pre-tokenized. Neural systems run the pipeline their
lemmatizer requires, including tagging. Every answer in a window is compared
with the system's answers in the accuracy run; udlex differed in
\speedMismatch{} of \speedChecked{} checked answers.

Table~\ref{tab:speed} shows the result. On one core, udlex's Rust core
processes a median of \speedRsOneMedian{} words per second
(\speedRsOneMin{} to \speedRsOneMax{} across languages) and its Python
binding \speedPyOneMedian. The nearest competitor is simplemma, which udlex
outruns by a factor of \speedsimplemmaUp{} on \speedsimplemmaN{} shared
languages; the only language where it does not win is Georgian (kat), where
simplemma is faster by a percent. On an RTX~4090, Stanza processes a median
of \speedstanzaGpu{} words per second, and udlex on one CPU core is still
\speedstanzaGpuOneUp{} times faster on the same languages. On 16
threads udlex reaches a median of \speedRsSixteenMedian{} words per second.
Opening a bundle and answering the first sentence takes a median of
\loadRsMedianMs{}~ms and at most \loadRsMaxMs~ms, for the largest bundle
(Finnish). udlex's peak memory includes the test text that the timing
process holds.

\paragraph{Measurement conditions.} The \cellsNew{} cells of udlex and of the
GPU systems were timed with the machine's other services paused, with a
median foreign load of \foreignNew{} cores. The \cellsReused{} CPU cells of
the other systems come from an earlier session on the same machine with the
same harness, when other processes used a median of \foreignReused{} cores;
their models and inputs do not depend on udlex's dictionaries. Each cell
records its foreign load.

\section{Related work}

Dictionary information already helps neural lemmatizers, as extra training
data \citep{bergmanis-goldwater-2019-data} or as candidate analyses at
inference \citep{milintsevich-sirts-2021-enhancing}, and lemma generation
without language-specific training has been studied with large language models
\citep{toporkov-etal-2025-lemma}. LemmaGen \citep{jursic-etal-2010-lemmagen}
learns ripple-down suffix rules, and UDPipe~1 \citep{straka-strakova-2017-tokenizing}
a tagger with lemma rules. udlex combines corpus tables, ordered dictionaries
and suffix rules and targets throughput and coverage rather than the last
point of accuracy.

\section{Conclusion}

udlex lemmatizes \relLangs{} languages from compiled dictionaries, is the
fastest measured lemmatizer on nearly every language, and is more accurate
than the other non-neural systems we compared. Neural toolkits remain the
choice when the last points of accuracy matter more than throughput. The
code is MIT-licensed; dictionaries, evaluation receipts and the scripts that
produced every table are public.

\section*{Limitations}

udlex lemmatizes one token at a time and cannot resolve forms whose lemma
depends on context; this is the main gap to Stanza and UDPipe~2. Two-fold
scores come from small treebanks whose halves may share documents. For the
languages without a development split, the dictionary builder introduced the
development carve-out after the first test evaluation, so those scores are
partly exploratory. Accuracy on languages without a treebank is unmeasured.
Speed numbers depend on the CPU and on the token mix of the UD test texts,
which repeat frequent words; the competitor CPU cells were timed under a
slightly higher background load than udlex's.

\section*{Ethical considerations}

The dictionaries are derived from public resources and keep their licenses
and attribution. Packages with noncommercial parts are marked; users must
follow each package's license. Lemmatization quality differs widely across
languages, which can disadvantage speakers of low-resource languages in
downstream systems.

\bibliography{references}

\end{document}
